\noindent\textbf{1. Lí do chọn đề tài}


Hiện nay, việc học tiếng Anh đối với trẻ em đang được nhiều phụ huynh quan tâm ngay từ khi còn nhỏ. Bên cạnh việc học trên lớp, nhiều gia đình còn cho trẻ tiếp cận thêm nhiều ứng dụng học tập trên điện thoại hoặc máy tính bảng để giúp việc học trở nên thú vị hơn.


Tuy nhiên, trên thực tế không phải ứng dụng nào cũng thực sự phù hợp với trẻ nhỏ. Có nhiều ứng dụng chứa quá nhiều lý thuyết, giao diện khó sử dụng hoặc thiếu tính tương tác khiến trẻ nhanh chán và không còn hứng thú học tập. Ngoài ra, một số ứng dụng còn chưa hỗ trợ tốt khả năng luyện nghe và phát âm cho trẻ em Việt Nam.


Trong quá trình tìm hiểu thực tế, em nhận thấy rằng trẻ em thường có xu hướng học tập hiệu quả hơn khi được tiếp cận bằng hình ảnh trực quan và sinh động, âm thanh thì sôi động và các hoạt động mang tính tương tác cao. Chính vì vậy, em lựa chọn phát triển đề tài ứng dụng học tiếng Anh dành cho trẻ em với mong muốn tạo ra một môi trường học tập thân thiện, dễ tiếp cận và tạo cảm giác hứng thú trong quá trình học.


Ứng dụng được xây dựng nhằm hỗ trợ trẻ làm quen với tiếng Anh theo hướng tự nhiên và chủ động hơn thông qua các chức năng như học từ vựng, luyện phát âm, trò chơi câu hỏi và một số tính năng ứng dụng trí tuệ nhân tạo. Ngoài việc hỗ trợ học tập cho trẻ, hệ thống còn giúp phụ huynh theo dõi kết quả và quá trình sử dụng ứng dụng một cách thuận tiện hơn.


\noindent\textbf{2. Đối tượng nghiên cứu}


Đối tượng nghiên cứu của đề tài là trẻ em trong độ tuổi từ 6 đến 10 tuổi. Đây là độ tuổi phù hợp để bắt đầu làm quen và phát triển khả năng học ngoại ngữ thông qua các hoạt động trực quan và tương tác.


\noindent\textbf{3. Phạm vi nghiên cứu}


Đề tài được thực hiện nhằm phát triển một ứng dụng hỗ trợ học tiếng Anh cho trẻ em trên nền tảng Android, đồng thời kết hợp với Website quản trị để quản lý dữ liệu và nội dung học tập của hệ thống.


\noindent\textbf{4. Kết cấu đề tài}


Nội dung kết cấu đề tài được chia thành các chương như sau:

\vspace{0.3cm}

\noindent\textbf{Chương 1: Cơ sở lý thuyết}

Giới thiệu những kiến thức và công nghệ được áp dụng trong quá trình phát triển ứng dụng.

\vspace{0.25cm}

\noindent\textbf{Chương 2: Khảo sát và Phân tích hệ thống}

Trình bày việc khảo sát nhu cầu người dùng và phân tích các chức năng hệ thống.

\vspace{0.25cm}

\noindent\textbf{Chương 3: Thiết kế hệ thống}

Mô tả cách tổ chức dữ liệu, giao diện và các thành phần chính của hệ thống.

\vspace{0.25cm}

\noindent\textbf{Chương 4: Cài đặt và triển khai hệ thống}

Trình bày quá trình xây dựng, kết nối và vận hành thử nghiệm ứng dụng.
